\section{Anexos}
\label{sec:anexos}

\subsection*{Anexo A: Instalación y configuración}
\label{anexo:instalacion}
Los resultados guardados se obtuvieron con Python 3.12.8 en Windows 11. Desde la raíz del repositorio se puede preparar el entorno con estos comandos; la instalación en un equipo limpio debe comprobarse por separado:
\begin{lstlisting}[language=bash, caption={Clonación, entorno virtual e instalación}]
git clone https://github.com/frankich99/PROYECTO_Algoritmos-Avanzados-UNSAAC.git
cd PROYECTO_Algoritmos-Avanzados-UNSAAC
python -m venv .venv
.\.venv\Scripts\Activate.ps1
python -m pip install -r python/requirements.txt
\end{lstlisting}
Los requisitos principales son Python 3.10 o superior, \texttt{pytest}, \texttt{networkx}, \texttt{osmnx}, \texttt{streamlit}, \texttt{plotly}, \texttt{pandas}, \texttt{numpy} y \texttt{matplotlib}. La compilación del informe requiere una distribución LaTeX con \texttt{pdflatex} y \texttt{biber}. La red OSM y el clima se versionan como \textit{snapshots}; no es obligatorio descargarlos para ejecutar las pruebas.

\subsection*{Anexo B: Uso del prototipo}
\label{anexo:uso}

\noindent\textbf{Suite de pruebas.} Desde la raíz:
\begin{lstlisting}[language=bash, caption={Ejecución de la suite completa de pruebas (278 pruebas)}]
.\.venv\Scripts\python.exe -m pytest python/tests -q
\end{lstlisting}

\noindent\textbf{Configuración de credenciales de API (opcional).} Tras revocar cualquier clave expuesta, configure una clave nueva y propia. Para la CLI se admite la variable de entorno, nunca un archivo versionado (\url{https://home.openweathermap.org/api_keys}):
\begin{lstlisting}[language=bash, caption={Definición de variable de entorno para OpenWeatherMap}]
# PowerShell en Windows
$env:OPENWEATHERMAP_API_KEY = "SU_CLAVE_NUEVA"
\end{lstlisting}
La web no toma automáticamente la variable ni archivos locales: exige introducir una clave nueva en el campo de contraseña de la consulta opcional. Si el proveedor rechaza la conexión, la interfaz conserva el archivo local sin reportar una nueva lectura. La clave no se imprime ni se guarda en el informe.

\noindent\textbf{Interfaz de comandos.} Desde \texttt{python/}:
\begin{lstlisting}[language=bash, caption={Comandos principales de la CLI}]
..\.venv\Scripts\python.exe main.py list
..\.venv\Scripts\python.exe main.py run --queue all --trace
..\.venv\Scripts\python.exe main.py compare
..\.venv\Scripts\python.exe main.py astar --input data/processed/cusco_osm_drive.json
\end{lstlisting}
El siguiente ejemplo recorre generación, validación, ejecución, exportación y verificación. Escribe dos archivos de ejemplo en \texttt{python/}, sin alterar los resultados guardados:
\begin{lstlisting}[language=bash, caption={Ciclo de prueba con una instancia generada}]
..\.venv\Scripts\python.exe main.py generate --type random --nodes 40 --seed 3 -o ejemplo.json
..\.venv\Scripts\python.exe main.py run -i ejemplo.json -q binary -o resultado.json
..\.venv\Scripts\python.exe main.py verify -i ejemplo.json -r resultado.json
..\.venv\Scripts\python.exe main.py compare -i ejemplo.json
\end{lstlisting}
\texttt{graph\_io.py} acepta JSON, TXT, CSV o aristas en línea; \texttt{run} exporta distancias, rutas y contadores en JSON. Una entrada con peso negativo se rechaza antes de ejecutar Dijkstra. Las mediciones comparativas se obtienen por separado con \texttt{reproduce\_canonical.py}.

\noindent\textbf{Aplicación web.} Desde \texttt{python/}:
\begin{lstlisting}[language=bash, caption={Ejecución de la aplicación Streamlit}]
..\.venv\Scripts\python.exe -m streamlit run app.py
\end{lstlisting}
La aplicación tiene cinco secciones visibles. En Ruta se seleccionan costo, recorrido, algoritmo y cola; se pueden contrastar las tres colas de Dijkstra y $A^*$ por costo, elegir una observación histórica, local o temporal como escenario uniforme, y descargar los nodos. El plano local no necesita Internet; solo el mapa base necesita teselas. Comparación filtra perfiles y escalas, muestra mediana e IQR y la relación frente al binario, admite un JSON temporal compatible y descarga CSV. Clima y tráfico separa observaciones y supuestos; Traza manual permite recorrer extracciones y comprobar el certificado; Datos y método presenta protocolo, procedencia, límites y descarga del informe si existe. Ninguna consulta opcional sobrescribe las mediciones.

\noindent\textbf{Reproducción de tablas y figuras.} Desde la raíz, sin alterar las mediciones del 28/10/2026:
\begin{lstlisting}[language=bash, caption={Regeneración de tablas y figuras con mediciones guardadas}]
.\.venv\Scripts\python.exe python/scripts/reproduce_canonical.py --skip-benchmark --skip-applied
\end{lstlisting}
El comando anterior reutiliza \texttt{benchmark\_results.json} y \texttt{applied\_results.json} y escribe tablas en \texttt{latex/tablas/} y figuras en \texttt{latex/figuras/}. Sin las opciones de omisión, el script repite las mediciones y \emph{sobrescribe} ambos JSON; ejecutar el protocolo completo solo en una copia separada. La red OSM y el clima ya almacenados no se descargan de nuevo.
El piloto del 04/11/2026 se conserva aparte en \texttt{pilot\_benchmark\_results.json}. Para regenerar únicamente su figura desde las muestras guardadas, ejecutar desde la raíz \texttt{python python/scripts/reproduce\_pilot.py --plot-only}. La imagen se guarda en \texttt{latex/figuras/piloto\_mediciones.png} y aparece en la Sección~\ref{sec:metodologia}.

\noindent\textbf{Compilación del informe.} Desde \texttt{latex/}:
\begin{lstlisting}[language=bash, caption={Compilación con pdflatex y biber}]
.\compilar.ps1
\end{lstlisting}

\subsection*{Anexo C: Cronología verificable del repositorio}
\label{anexo:bitacora}
La Tabla~\ref{tab:bitacora_detallada} reúne cambios del historial Git hasta el 04/11/2026. Cada fecha y autor corresponden al commit, no necesariamente al momento de trabajo o a la persona que redactó cada línea. No se deducen reuniones, horas de dedicación ni lecturas individuales. La bitácora de aprendizaje por estudiante sigue sin constancia en el repositorio.

\begin{table}[H]
    \centering
    \fontsize{8.5}{10.5}\selectfont
    \renewcommand{\arraystretch}{1.12}
    \setlength{\tabcolsep}{3.5pt}
    \caption{Cambios fechados comprobables en Git}
    \label{tab:bitacora_detallada}
    \begin{tabularx}{\textwidth}{@{}>{\bfseries\raggedright\arraybackslash}p{2.65cm} >{\centering\arraybackslash}p{1.55cm} >{\raggedright\arraybackslash}p{1.35cm} >{\raggedright\arraybackslash}X@{}}
        \toprule
        \textbf{Autor en Git} & \textbf{Fecha} & \textbf{Commit} & \textbf{Cambios identificables} \\
        \midrule
        Choquenaira, Noe & 02/10/2026 & \texttt{592589a} & \texttt{dijkstra.py}, \texttt{indexed\_min\_heap.py}, \texttt{cost\_model.py} y primera versión del informe. \\
        \midrule
        Porroa, Yeni & 04/10/2026 & \texttt{77983b7} & \texttt{d\_ary\_heap.py}, capítulo de evaluación del stack y tabla comparativa. \\
        \midrule
        Quispe, Romario & 04/10/2026 & \texttt{72df976} & \texttt{fibonacci\_heap.py}, \texttt{scalability\_benchmark.py} y resultados preliminares. \\
        \midrule
        Yaranga, Aldo & 04/10/2026 & \texttt{d50e67b} & \texttt{test\_dijkstra.py}, cuaderno Colab y exportación de figuras y tablas. \\
        \midrule
        Choquenaira, Noe & 23/10/2026 & \texttt{b307be4} & \texttt{a\_star.py} y capítulo de aprendizaje. \\
        Porroa, Yeni & 23/10/2026 & \texttt{3b91a5e} & Capítulo de diseño algorítmico y análisis de estructuras. \\
        Quispe, Romario & 23/10/2026 & \texttt{e3af921} & Actualización de \texttt{benchmark\_results.json}; el snapshot canónico se integró después como commit grupal. \\
        Yaranga, Aldo & 23/10/2026 & \texttt{65d303d} & \texttt{cli.py} y \texttt{test\_cli.py}. \\
        \midrule
        Grupo 5 & 28/10/2026 & \texttt{7922f52} & \textit{Snapshots} OSM y clima, banco de 1,350 tiempos y caso aplicado. \\
        \midrule
        Grupo 5 & 03/11/2026 & \texttt{4025f8f} & Revisión de estructura e índice del informe, PDF compilado. \\
        \midrule
        Grupo 5 & 04/11/2026 & \texttt{ced0b82} & Ajustes metodológicos y vista de referencia OSM. \\
        \midrule
        Grupo 5 & 04/11/2026 & \texttt{c119aa6} & Piloto reproducible, figura y síntesis comparativa revisada. \\
        \bottomrule
    \end{tabularx}
\end{table}

\begin{table}[H]
    \centering
    \fontsize{9}{11}\selectfont
    \renewcommand{\arraystretch}{1.15}
    \caption{Anexo D: matriz de trazabilidad entre objetivos, componentes y evidencia}
    \label{tab:matriz_trazabilidad}
    \begin{tabularx}{\textwidth}{@{}>{\bfseries\raggedright\arraybackslash}p{2.6cm} >{\raggedright\arraybackslash}p{4.7cm} >{\raggedright\arraybackslash}X@{}}
        \toprule
        \textbf{Objetivo} & \textbf{Componente} & \textbf{Evidencia} \\
        \midrule
        OE1 & Modelo de grafo y costo temporal & Sección~\ref{sec:formulacion}; ecuación~\eqref{eq:costo}; instancia ilustrativa y red OSM. \\
        OE2 & Tres colas y Dijkstra parametrizado & Sección~\ref{sec:diseno_algoritmico}; \texttt{src/datastructures/}; 61 pruebas de contrato. \\
        OE3 & Validación y certificado & Sección~\ref{sec:pruebas}; 278 pruebas; Bellman y cadenas al origen; NetworkX y Bellman-Ford. \\
        OE4 & Benchmark, caso OSM, $A^*$ y clima & Sección~\ref{sec:resultados}; 1350 muestras; 891 vértices; 823 vs 887 expansiones; factor 1.0691. \\
        AG\_C06.01 & Aprendizaje autónomo & Sección~\ref{sec:fundamentos}; Anexo~C; pendientes individuales en \texttt{docs/CONTROL\_ENTREGA.md}. \\
        AG\_C06.02 & Selección de stack & Sección~\ref{sec:stack}; Tabla~\ref{tab:matriz_multicriterio}; adopción temprana por verificar. \\
        AG\_C06.03 & Experimentación y pensamiento crítico & Sección~\ref{sec:metodologia}; Sección~\ref{sec:discusion}; conclusión condicionada. \\
        Aplicación & Streamlit integrado & Sección~\ref{sec:implementacion}; \texttt{app.py}; pruebas de humo. \\
        \bottomrule
    \end{tabularx}
\end{table}

\subsection*{Anexo E: Matriz de contribuciones}
\label{anexo:contribuciones}
\begin{table}[H]
    \centering
    \fontsize{9}{11}\selectfont
    \renewcommand{\arraystretch}{1.15}
    \caption{Contribuciones por integrante y evidencia en el repositorio}
    \label{tab:contribuciones}
    \begin{tabularx}{\textwidth}{@{}>{\bfseries\raggedright\arraybackslash}p{2.8cm} >{\raggedright\arraybackslash}p{5.0cm} >{\raggedright\arraybackslash}X@{}}
        \toprule
        \textbf{Integrante} & \textbf{Componentes principales} & \textbf{Evidencia} \\
        \midrule
        Choquenaira, Noe & Dijkstra, heap binario, modelo didáctico, $A^*$ e informe & \texttt{592589a}, \texttt{b307be4}, \texttt{f0799bf}; archivos en anexo C. \\
        Porroa, Yeni & Heap 4-ario y análisis de estructuras & \texttt{77983b7}, \texttt{3b91a5e}; \texttt{d\_ary\_heap.py} y Sección~\ref{sec:diseno_algoritmico}. \\
        Quispe, Romario & Fibonacci heap, benchmark y resultados preliminares & \texttt{72df976}, \texttt{e3af921}; módulos y JSON del anexo C. \\
        Yaranga, Aldo & Pruebas, CLI, cuaderno y exportadores & \texttt{d50e67b}, \texttt{65d303d}; pruebas y CLI del anexo C. \\
        \bottomrule
    \end{tabularx}
\end{table}

\subsection*{Anexo F: Estado del seguimiento}
\label{anexo:seguimiento}
Al 04/11/2026 hay código integrado, pruebas automatizadas, análisis de complejidad, datos del banco y de la consulta OSM, tablas y figuras, instrucciones de reproducción y un PDF del informe. El piloto de 18 muestras se guardó el 04/11, \emph{después} de las mediciones canónicas del 28/10; no se presenta como una etapa previa que no ocurrió. Las muestras cronometradas de los JSON se expresan en microsegundos; los 387.197 segundos del caso OSM corresponden al costo estimado de viaje de la ruta.

No se localizaron bitácoras originales de fuentes, actividades y aprendizajes por integrante, ni actas con compromisos fechados. Tampoco hay constancia de recepción de la segunda entrega o de una presentación ante el docente. El equipo debe confirmar la instalación en un entorno limpio, conservar una salida reciente de la suite de pruebas y aportar sus registros individuales auténticos si existen. El 04/11 es una fecha objetivo comunicada por el equipo, no una fecha oficial ni una entrega acreditada en el repositorio. La matriz de pendientes y la plantilla están en \texttt{docs/CONTROL\_ENTREGA.md} y \texttt{docs/plantilla\_evidencia\_individual.md}; no sustituyen esas evidencias.
